\documentclass[14pt]{beamer}
\usepackage{pepslides}
\title{Plane wave propagation}
\lecturenumber{16}
\hypersetup{pdftitle={Lecture 16: Plane wave propagation}}

\begin{document}

\titleslide

\begin{frame}[t]
\heading{From last time\ldots}
Linearised equations of motion including a stress glut:
\begin{eqn}
\begin{gathered}
\rho\frac{\partial^2 u_{i}}{\partial t^2} - \frac{\partial}{\partial x_{j}}\left(A_{ijkl}\frac{\partial u_{k}}{\partial x_{l}}\right) = -\frac{\partial\overline{S}_{ij}}{\partial x_{j}},\\[\eqgap]
A_{ijkl}\frac{\partial u_{k}}{\partial x_{l}}\hat{n}_{j} = \overline{S}_{ij}\hat{n}_{j}.
\end{gathered}
\end{eqn}
Definition of the elastic tensor:
\begin{eqn}
A_{ijkl}(\mathbf{x}) = \frac{\partial^2 W}{\partial F_{ij}\partial F_{kl}}(\mathbf{x}, \mathbf{I}).
\end{eqn}
\end{frame}

\begin{frame}
\centering
\includegraphics[height=0.9\paperheight]{vandriel_nissenmeyer_global_wavefield.png}
\credit{Martin van Driel \& Tarje Nissen-Meyer}
\end{frame}

\sectionslide{The Christoffel equation}

\begin{frame}
We consider an infinite and homogeneous body.

Look for a plane wave solution of the form:
\begin{eqn}
u_{i}(\mathbf{x},t)=f(t-p_{j}x_{j})a_{i}.
\end{eqn}
The \textbf{slowness vector}, $\mathbf{p}$, sets the orientation and speed of the wavefronts.

The \textbf{polarisation vector}, $\mathbf{a}$, sets the direction of particle motion.

The \textbf{pulse shape} is set by $f$, and does not vary as the wave propagates -- in this
sense the waves are \textbf{non-dispersive} (but\ldots).
\end{frame}

\begin{frame}
\begin{columns}[c]
\begin{column}{0.45\textwidth}
Wavefronts:
\begin{eqn}
t-p_{j}x_{j}=\mathrm{constant}.
\end{eqn}
Phase speed:
\begin{eqn}
c=\frac{1}{\|\mathbf{p}\|}.
\end{eqn}
\end{column}
\begin{column}{0.5\textwidth}
\centering
\begin{tikzpicture}[rotate=35,scale=1.55,>={Stealth[length=13pt,width=10pt]}]
\foreach \x in {0,0.75,...,3}
  \draw[line width=2.4pt,dash pattern=on 5pt off 4pt] (\x,-3.3) -- (\x,1.7);
\draw[->,line width=2.4pt] (3,0.8) -- ++(1.3,0) node[right=3pt] {\eqsize$\mathbf{p}$};
\draw[->,line width=2.4pt,red!85!black] (3,0.8) -- ++(62:1.9) node[above=3pt] {\eqsize$\mathbf{a}$};
\end{tikzpicture}
\end{column}
\end{columns}
\end{frame}

\begin{frame}
For the plane wave ansatz:
\begin{eqn}
\begin{gathered}
\frac{\partial^{2}u_{i}}{\partial t^{2}}=f''(t-p_{j}x_{j})a_{i},\\[\eqgap]
\frac{\partial}{\partial x_{j}}\left(A_{ijkl}\frac{\partial u_{k}}{\partial x_{l}}\right)=f''(t-p_{m}x_{m})A_{ijkl}p_{j}p_{l}a_{k}.
\end{gathered}
\end{eqn}
The equations of motion are satisfied if and only if:
\begin{eqn}
(\rho\delta_{ik}-A_{ijkl}p_{j}p_{l})a_{k}=0.
\end{eqn}
This is known as the \textbf{Christoffel equation} -- it places a constraint on both the
slowness and polarisation vectors.
\end{frame}

\begin{frame}
It will be useful to define the \textbf{Christoffel matrix}, which is a function of $\mathbf{p}$:
\begin{eqn}
\Gamma_{ik}(\mathbf{p})=\frac{1}{\rho}A_{ijkl}p_{j}p_{l}.
\end{eqn}
The Christoffel equation can then be written:
\begin{eqn}
[\bm{\Gamma}(\mathbf{p})-\mathbf{I}]\mathbf{a}=\mathbf{0}.
\end{eqn}
This is a \textbf{homogeneous linear equation} for $\mathbf{a}$, that can have non-trivial solutions
if and only if:
\begin{eqn}
\det[\bm{\Gamma}(\mathbf{p})-\mathbf{I}]=0.
\end{eqn}
The roots of this equation define the \textbf{slowness surface}.
\end{frame}

\sectionslide{Reduction to an eigenvalue problem}

\begin{frame}
From the definition of the Christoffel matrix, we have:
\begin{eqn}
\bm{\Gamma}(\gamma \mathbf{p})=\gamma^{2}\bm{\Gamma}(\mathbf{p}).
\end{eqn}
It is a \textbf{homogeneous function of order two}.

The hyperelastic symmetry implies that the Christoffel matrix is \textbf{symmetric}:
\begin{eqn}
\Gamma_{ki}(\mathbf{p})=\frac{1}{\rho}A_{kjil}p_{j}p_{l}=\frac{1}{\rho}A_{ilkj}p_{j}p_{l}=\frac{1}{\rho}A_{ijkl}p_{j}p_{l}=\Gamma_{ik}(\mathbf{p}).
\end{eqn}
\end{frame}

\begin{frame}
Fixing the orientation of the slowness vector, we can write:
\begin{eqn}
\mathbf{p}=\frac{1}{c}\hat{\mathbf{p}},
\end{eqn}
and hence the Christoffel equation reduces to:
\begin{eqn}
\bm{\Gamma}(\hat{\mathbf{p}})\mathbf{a}=c^{2}\mathbf{a}.
\end{eqn}
This is a \textbf{symmetric eigenvalue problem}:
\begin{enumerate}
\item Squared phase speeds are \textbf{real}, generally \textbf{three distinct values}
\item Polarisation vectors for distinct squared phase speeds are \textbf{orthogonal}
\item Degenerate squared phase speeds are possible
\end{enumerate}
\end{frame}

\begin{frame}
The slowness surface has three sheets, which can touch along some directions:
\begin{center}
\includegraphics[width=0.74\textwidth]{L4F1.pdf}
\end{center}
\vspace{-0.6em}
\begin{eqn}
\begin{gathered}
{}[\|\mathbf{p}\|^{2}c_{1}(\hat{\mathbf{p}})^{2}-1][\|\mathbf{p}\|^{2}c_{2}(\hat{\mathbf{p}})^{2}-1][\|\mathbf{p}\|^{2}c_{3}(\hat{\mathbf{p}})^{2}-1]=0,\\[\eqgap]
c_{1}(\hat{\mathbf{p}})\le c_{2}(\hat{\mathbf{p}})\le c_{3}(\hat{\mathbf{p}}).
\end{gathered}
\end{eqn}
\end{frame}

\begin{frame}
\centering
\animation[width=0.46\textwidth]{anisotropy_iso.pdf}{anisotropy_iso.webm}\hspace{1.5em}%
\animation[width=0.46\textwidth]{anisotropy_ti.pdf}{anisotropy_ti.webm}

\raggedright
A point force in an isotropic medium (left) and in a strongly anisotropic one (right). The quasi-P and quasi-S
wavefronts are not circles, and the quasi-S front has cusps. These are two-dimensional solutions, with motion in
the plane, so only two of the three sheets of the slowness surface appear.
\end{frame}

\sectionslide{Isotropic media}

\begin{frame}
We have seen that the elastic tensor in an isotropic solid takes the form:
\begin{eqn}
A_{ijkl}=\lambda\delta_{ij}\delta_{kl}+\mu(\delta_{ik}\delta_{jl}+\delta_{il}\delta_{jk}).
\end{eqn}
Putting into the definition of the Christoffel matrix, we arrive at:
\begin{eqn}
\begin{aligned}
\Gamma_{ik}(\hat{\mathbf{p}})&=\frac{\lambda+\mu}{\rho}\hat{p}_{i}\hat{p}_{k}+\frac{\mu}{\rho}\delta_{ik}\\[\eqgap]
&=\frac{\lambda+2\mu}{\rho}\hat{p}_{i}\hat{p}_{k}+\frac{\mu}{\rho}(\delta_{ik}-\hat{p}_{i}\hat{p}_{k}).
\end{aligned}
\end{eqn}
In the final expression, note the \textbf{orthogonal projections} parallel and perpendicular
to the direction of the slowness vector.
\end{frame}

\begin{frame}
Isotropic Christoffel matrix:
\begin{eqn}
\Gamma_{ik}(\hat{\mathbf{p}})=\frac{\lambda+2\mu}{\rho}\hat{p}_{i}\hat{p}_{k}+\frac{\mu}{\rho}(\delta_{ik}-\hat{p}_{i}\hat{p}_{k}).
\end{eqn}
By inspection we identify an eigenvector and eigenvalue:
\begin{eqn}
\mathbf{a}\propto\hat{\mathbf{p}},\qquad \alpha=\sqrt{\frac{\lambda+2\mu}{\rho}}.
\end{eqn}
These plane waves are known as \textbf{P-waves}:
\begin{itemize}
\item Particle motion along the propagation direction
\item Phase speed the same in all directions
\item Associated with a non-degenerate eigenspace
\end{itemize}
\end{frame}

\begin{frame}
Isotropic Christoffel matrix:
\begin{eqn}
\Gamma_{ik}(\hat{\mathbf{p}})=\frac{\lambda+2\mu}{\rho}\hat{p}_{i}\hat{p}_{k}+\frac{\mu}{\rho}(\delta_{ik}-\hat{p}_{i}\hat{p}_{k}).
\end{eqn}
By inspection we identify eigenvectors and an eigenvalue:
\begin{eqn}
\mathbf{a}\perp\hat{\mathbf{p}},\qquad \beta=\sqrt{\frac{\mu}{\rho}}.
\end{eqn}
These plane waves are known as \textbf{S-waves}:
\begin{itemize}
\item Particle motion perpendicular to the propagation direction
\item Phase speed the same in all directions
\item Associated with a two-fold degenerate eigenspace
\end{itemize}
\end{frame}

\sectionslide{Slightly anisotropic media}

\begin{frame}[t]
\heading{Anisotropy in the Earth}
\centering
\includegraphics[height=0.68\paperheight]{anisotropy_rocks_photos.png}
\end{frame}

\begin{frame}
We consider a slightly anisotropic elastic tensor:
\begin{eqn}
A_{ijkl}=\lambda\delta_{ij}\delta_{kl}+\mu(\delta_{ik}\delta_{jl}+\delta_{il}\delta_{jk})+s\,\delta A_{ijkl}.
\end{eqn}
Resulting Christoffel matrix:
\begin{eqn}
\Gamma_{ik}(\hat{\mathbf{p}})=\Gamma^{0}_{ik}(\hat{\mathbf{p}})+s\,\delta\Gamma_{ik}(\hat{\mathbf{p}}),\qquad
\delta\Gamma_{ik}(\hat{\mathbf{p}})=\frac{1}{\rho}\delta A_{ijkl}\hat{p}_{j}\hat{p}_{l}.
\end{eqn}
Eigenvalue problem to be solved:
\begin{eqn}
[\bm{\Gamma}^{0}(\hat{\mathbf{p}})+s\,\delta\bm{\Gamma}(\hat{\mathbf{p}})]\mathbf{a}=c^{2}\mathbf{a}.
\end{eqn}
\end{frame}

\begin{frame}
Apply \textbf{non-degenerate} perturbation theory:
\begin{eqn}
\begin{gathered}
c=\alpha+s\,\delta\alpha+\cdots,\\
\mathbf{a}=\hat{\mathbf{p}}+s\,\delta \mathbf{a}+\cdots.
\end{gathered}
\end{eqn}
First-order term in eigenvalue problem:
\begin{eqn}
\bm{\Gamma}^{0}(\hat{\mathbf{p}})\delta \mathbf{a}+\delta\bm{\Gamma}(\hat{\mathbf{p}})\hat{\mathbf{p}}=\alpha^{2}\delta\mathbf{a}+2\alpha\,\delta\alpha\,\hat{\mathbf{p}}.
\end{eqn}
Through a standard argument, we arrive at:
\begin{eqn}
\delta\alpha=\frac{1}{2\alpha}\hat{\mathbf{p}}\cdot\delta\bm{\Gamma}(\hat{\mathbf{p}})\hat{\mathbf{p}}.
\end{eqn}
Polarisation \textbf{almost aligned} with propagation direction, phase speed now \textbf{varies
with the orientation}. This is a \textbf{quasi P-wave}.
\end{frame}

\begin{frame}
Apply \textbf{degenerate} perturbation theory:
\begin{eqn}
\begin{gathered}
c=\beta+s\,\delta\beta+\cdots,\\
\mathbf{a}=a_{1}\hat{\mathbf{t}}_{1}+a_{2}\hat{\mathbf{t}}_{2}+s\,\delta \mathbf{a}+\cdots.
\end{gathered}
\end{eqn}
Leads to a $2\times2$ symmetric eigenvalue problem:
\begin{eqn}
\begin{pmatrix} \hat{\mathbf{t}}_{1}\cdot\delta\bm{\Gamma}(\hat{\mathbf{p}})\hat{\mathbf{t}}_{1} & \hat{\mathbf{t}}_{1}
    \cdot\delta\bm{\Gamma}(\hat{\mathbf{p}})\hat{\mathbf{t}}_{2} \\ \hat{\mathbf{t}}_{2}\cdot\delta\bm{\Gamma}(\hat{\mathbf{p}})\hat{\mathbf{t}}_{1} &
    \hat{\mathbf{t}}_{2}\cdot\delta\bm{\Gamma}(\hat{\mathbf{p}})\hat{\mathbf{t}}_{2} \end{pmatrix} \begin{pmatrix} a_{1}
         \\ a_{2} \end{pmatrix} = 2\beta\,\delta\beta \begin{pmatrix} a_{1} \\ a_{2} \end{pmatrix}.
\end{eqn}
Solution gives \textbf{two perturbed phase speeds} which are usually distinct -- the
initial \textbf{degeneracy is split}.

Zeroth-order eigenvectors orthogonal to each other and to the propagation direction; the
first-order correction will have a component along the propagation direction.

These are \textbf{quasi S-waves}, and this behaviour is known as \textbf{shear wave splitting}.
\end{frame}

\begin{frame}
\centering
\includegraphics[width=0.85\textwidth]{long2009_splitting_flow.png}
\credit{Long \& Silver (2009)}

\raggedright
Convection in the mantle creates large-scale \textbf{alignment} of anisotropic mineral
grains with the \textbf{flow direction} -- produces an effective seismic anisotropy.
Example here shows the relation between shallow anisotropy and plate motion.
\end{frame}

\begin{frame}
The \textbf{interaction of plane waves with boundaries} is not treated quantitatively in this course.

The usual \textbf{reflections} and \textbf{transmissions} of waves occur when a plane elastic
wave interacts with a boundary.

The only interesting feature is the \textbf{conversion} between different wave types. For
example, in an isotropic medium, an incident P-wave can produce reflected and
transmitted P- and S-waves.
\end{frame}

\begin{frame}
\centering
\animation[height=0.72\paperheight]{boundary_t1.pdf}{boundary.webm}

\raggedright
A P-wave from a point source reaching a boundary with a faster medium.
\end{frame}

\begin{frame}
\centering
\animation[height=0.72\paperheight]{boundary_t2.pdf}{boundary.webm}

\raggedright
Reflected and transmitted P- and S-waves, and head waves along the boundary. Arrows show the particle motion.
\end{frame}

\begin{frame}[t]
\heading{What you need to know and be able to do}
\begin{enumerate}
\item[(i)] How to derive the Christoffel equation from the linearised equations of motion.
\item[(ii)] Reduction of the problem to a symmetric eigenvalue problem, and a qualitative description of its solution.
\item[(iii)] How to solve the Christoffel equation for an isotropic solid, and to deduce the properties of plane P-
and S-waves.
\item[(iv)] How to use first-order perturbation theory to study the effect of slight anisotropy on plane waves.
\end{enumerate}
\end{frame}

\end{document}
